\chapter{Next Steps}
\label{ch:nextSteps}

The project has several areas that could be investigated in more detail in the future.
Most notably there was not an opportunity to use the radar system on a mountainside with snowpack.
Testing was limited to simple cases with a clear target (corner reflector) in distances with no snow presence.
Using the system to attempt avalanche detection will require many improvements.
These improvements would include autonomous data collection and storage,
a weatherproof enclosure, and more powerful data processing tools to quantitatively detect avalanches.


Another area to be investigated further is hardware.
Two B210s can be wired and synchronized together, effectively increasing the number of channels by two enabling more sophisticated beamforming.
Alternatively, more advanced and expensive SDRs could be considered.
Further hardware could be added to the front end of the devices to increase transmission power and improve signal strength and detection range or operational frequency.


Data visualization could be improved by real-time processing and displaying collected radar samples as a waterfall plot.
Systems could also be designed to automatically detect and manage data by calculating the detected power and determining if noticeable motion occurred.
That data could be saved while data with little to no motion could be discarded.
This would allow the system to be more autonomous only collecting and saving relevant information instead of a system that quickly fills up storage space.

